%% ============================================================
%%  SECCIÓN 2 — Resumen
%% ============================================================

\section{Resumen}
\label{sec:resumen}

En el presente informe de laboratorio se determinaron y validaron experimentalmente las capacidades metrológicas del osciloscopio digital (Rohde \& Schwarz HMO1202 / RTC1002 de \qty{200}{\mega\hertz}) en la medición de voltajes continuos (DC), tensiones alternas periódicas (AC) y en su funcionamiento como graficador paramétrico $X$-$Y$.

En primer lugar, para voltajes continuos de una pila y una fuente de alimentación de múltiples salidas (\qty{1,5}{\volt}, \qty{2,5}{\volt}, \qty{4,5}{\volt}, \qty{7,5}{\volt} y \qty{10,5}{\volt}), se comprobó que la precisión de lectura depende fuertemente de la escala vertical seleccionada. En escalas óptimas (máxima deflexión vertical sin saturación), la correlación lineal entre la lectura del osciloscopio $V_{\text{osc}}$ y el multímetro patrón $V_{\text{mult}}$ arrojó una pendiente de ajuste $m = 1{,}0036 \pm 0{,}0039$, intercepto $b = -0{,}021 \pm 0{,}026\,\unit{\volt}$ con un coeficiente de determinación $R^2 = 0{,}99995$, y errores relativos porcentuales contenidos entre \qty{0,12}{\percent} y \qty{0,45}{\percent}.

En segundo lugar, para la señal alterna sinusoidal de un transformador de \qty{6}{\volt} a \qty{60}{\hertz}, se determinó un voltaje pico a pico experimental de $V_{pp} = \qty{19,50 \pm 0,50}{\volt}$ (frente al patrón teórico de \qty{19,32}{\volt}, error relativo de \qty{0,93}{\percent}) y un voltaje eficaz calculado de $V_{\text{rms}} = \qty{6,79 \pm 0,14}{\volt}$ (frente a \qty{6,83}{\volt} en multímetro True-RMS, error de \qty{0,61}{\percent}). El período medido a \qty{2}{\milli\second\per\division} fue de $T = \qty{16,60 \pm 0,20}{\milli\second}$, equivalente a una frecuencia de $f = \qty{60,24 \pm 0,73}{\hertz}$ ($e_r = \qty{0,40}{\percent}$).

En tercer lugar, con el generador de funciones OWON AG 051F a \qty{250}{\hertz} y amplitud $V_p = \qty{4,80 \pm 0,20}{\volt}$, se midieron períodos de $T = \qty{4,00 \pm 0,05}{\milli\second}$ ($f = \qty{250,0 \pm 3,1}{\hertz}$) para ondas senoidales, triangulares y cuadradas, encontrándose el valor nominal plenamente contenido dentro del intervalo de incertidumbre experimental.

Finalmente, en modo $X$-$Y$ se sintetizaron figuras de Lissajous estables para relaciones de frecuencia $1:1$, $1:2$, $1:3$ y $1:4$, verificando la reflexión geométrica vertical al invertir el canal $\text{CH2}$ ($y \to -y$) y la proyección de un punto estático en $(+1{,}5\,\unit{\volt}, +3{,}0\,\unit{\volt})$ para señales continuas DC.
